\documentstyle[%


righttag%


,11pt%


,amssymb%


,russian%               remove in Eng.


,izvru%                 Set izven% in Eng.


% ,izvru-ks             for brief communications


]{amsart}





%      Math definitions





\newcommand{\la}{\langle}


\newcommand{\ra}{\rangle}


\newcommand{\up}[2]{\overset{#1}{#2}{}}


\newcommand{\tts}[1]{}


\renewcommand{\baselinestretch}{0.96}





%\hoffset=-15mm%                  For Eng.





\setcounter{page}{21}


\god{1998}


\nomer{9~(436)} %                        Remove in Eng.


%\nomer{Vol.\,42, No.\,9}%               Set in Eng.


%\str{\pageref{begin}--\pageref{end}}%  Set in Eng.


\udk{518\,:\,517.948}


\author{\it Ю.Г.\,Булычев}


% NB:   ^^ remove in Eng.


\title{Построение параметризованных решений линейных\\


операторных уравнений на основе модифицированного\\


метода Галеркина}





\date{}


% \pagestyle{myheadings}%         Set in Eng.





\pagestyle{plain} %              Remove in Eng.


\begin{document}





% \thispagestyle{plain}%          Set in Eng.





\maketitle


\label{begin}





{\bf 1. Введение.} При решении широкого класса прикладных задач


возникает необходимость построения так называемых параметризованных


(известных с точностью до параметров) аналитических решений,


которые с допустимой погрешностью удовлетворяют рассматриваемым


операторным уравнениям для заданных областей возможных значений


характерных параметров исследуемых задач.





Ниже излагается единый подход к построению параметризованных


аналитических решений линейных операторных уравнений (алгебраических,


дифференциальных, интегральных и т.д.), в основу которого положена


непрерывная зависимость решений от параметров и впервые


опубликованная в статьях [1], [2], [3] идея использования опорных


решений. Развитие теории осуществляется в новой


проективно-параметрической


постановке, являющейся обобщением теории проекционных методов


типа Галеркина [4]--[7] на рассматриваемый автором случай.


\medskip





{\bf 2. Проективно-параметрическая постановка задачи.} Пусть имеются


два нормированных пространства $X$ и $Y$ и в каждом из них выделено


полное подпространство $\wt X\subset X$ и $\wt Y\subset Y$. Будем


предполагать, что имеется непрерывный линейный оператор


$\wt\Phi(\omega)$, проектирующий


$Y$ на $\wt Y$, а также линейное операторное уравнение (которое по


аналогии с


[4]


будем называть точным)


\begin{equation}


K_1(\omega)x(\omega)\equiv G(\omega)x(\omega)-


\lambda T(\omega)x(\omega)=y(\omega),\quad x\in X,\ \; y\in Y,


\end{equation}


где $G(\omega)$ и $T(\omega)$ --- непрерывные линейные


параметризованные операторы,


отображающие $X$ в $Y$ и зависящие от вещественного векторного


параметра $\omega\in\Omega\subset R^M$ ($\Omega$ --- открытая


ограниченная связная область, $M\in\{1,2,\dotsc\}$),


$\lambda$ --- константа.





Кроме того, относительно оператора $G(\omega)$ для всех


$\omega\in\Omega$ будем предполагать


\begin{itemize}


\item[1)] оператор $G(\omega)$ имеет непрерывный обратный;


\item[2)] оператор $G(\omega)$ устанавливает взаимно однозначное


соответствие


между $\wt X$ и $\wt Y$, т.\,е. $G(\omega)\wt X=\wt Y$,


и, следовательно,


$G^{-1}(\omega)\wt Y=\wt X$.


\end{itemize}





Наряду с (1) рассмотрим приближенное уравнение


\begin{equation}


\wt K_1(\omega)\wt x(\omega)\equiv G(\omega)\wt x(\omega)-


\lambda\wt T(\omega)\wt x(\omega)=\wt\Phi(\omega)y(\omega),


\end{equation}


где $\wt T(\omega)$ --- непрерывный линейный параметризованный оператор


из


$\wt X$ в $\wt Y$.





Введем для всех $\omega\in\Omega$ следующие условия:


\begin{itemize}


\item[1)] для всякого $\wt x(\omega)\in\wt X$ имеет место неравенство


\begin{equation}


\|\wt\Phi(\omega)T(\omega)\wt x(\omega)-\wt T(\omega)


\wt x(\omega)\|\le\mu_0\|\wt x(\omega)\|,


\end{equation}


\item[2)] для каждого $x(\omega)\in X$ найдется такое


$\wt y_1(\omega)\in\wt Y$, что


\begin{equation}


\|T(\omega)x(\omega)-\wt y_1(\omega)\|\le\mu_1\|x(\omega)\|,


\end{equation}


\item[3)] существует такой элемент $\wt y_2(\omega)\in\wt Y$, что


\begin{equation}


\|y(\omega)-\wt y_2(\omega)\|\le\mu_2\|y_1(\omega)\|.


\end{equation}


\end{itemize}


Несложно показать, что при сделанных предположениях изучение


уравнений (1) и (2) сводится к изучению следующих уравнений ---


точного


\begin{equation}


K(\omega)x(\omega)\equiv G^{-1}(\omega)K_1(\omega)x(\omega)


\equiv x(\omega)-\lambda G^{-1}(\omega)T(\omega)x(\omega)=


G^{-1}(\omega)y(\omega)


\end{equation}


и соответствующего ему приближенного


\begin{equation}


\wt K(\omega)\wt x(\omega)\equiv G^{-1}(\omega)\wt K_1(\omega)


\wt x(\omega)


\equiv \wt x(\omega)-\lambda G^{-1}(\omega)\wt T(\omega)\wt x(\omega)=


G^{-1}(\omega)\wt\Phi(\omega)y(\omega).


\end{equation}





Пусть теперь имеется последовательность приближенных уравнений


и полученных с их помощью приближенных решений. Введем следующие


обозначения


$\wt X=X_n$, $\wt K(\omega)=K_n(\omega)$,


$\wt T(\omega)=T_n(\omega)$,


$\wt\Phi(\omega)=\Phi_n(\omega)$, где $n=1,2,\dotsc$


В этом случае вместо (6) и (7) можно ограничиться


более компактной и удобной на практике формой представления точного


и приближенного уравнений --- точное


\begin{equation}


K(\omega)x(\omega)\equiv x(\omega)-\lambda F(\omega)


x(\omega)=f(\omega),\quad x\in X,


\end{equation}


и приближенное


\begin{equation}


K_n(\omega)x_n(\omega)\equiv x_n(\omega)-


\lambda F_n(\omega)x_n(\omega)=P_n(\omega)f(\omega),\quad


x_n\in X,


\end{equation}


где


\begin{alignat*}{2}


F(\omega)&=G^{-1}(\omega)T(\omega), &\quad


f(\omega)&=G^{-1}(\omega)y_1(\omega), \\


F_n(\omega)&=G^{-1}(\omega)T_n(\omega), &\quad


P_n(\omega)&=G^{-1}(\omega)\Phi_n(\omega)G(\omega).


\end{alignat*}





Зададим теперь последовательность $\{Z_n\}$ полных пространств


$Z_n$, которые в частном случае могут совпадать с замкнутыми


подпространствами $X_n$ пространства $X$.


Будем считать, что в полном пространстве $X_n\subset X$ определен


непрерывный линейный параметризованный оператор $H_n(\omega)$,


отображающий


$X_n$ взаимно однозначно на полное пространство $Z_n$. В силу его


непрерывности и полноты следует непрерывность обратного оператора


$H^{-1}_n(\omega)$ для всех $\omega\in\Omega$. Полагаем также


существование непрерывного


линейного оператора $\wt H_n(\omega)$,


являющегося продолжением оператора


$H_n(\omega)$ на все


$X$, т.\,е. $\wt H_n(\omega)$ отображает $X$ на $Z_n$ и совпадает с


$H_n(\omega)$ на $X_n$ для всех $\omega\in\Omega$.





Пусть непрерывный линейный оператор $P_n(\omega)$ проектирует $X$


на $X_n$, при этом $P^2(\omega)X=X_n$, $P^2_n(\omega)=P_n(\omega)$,


где


$\omega\in\Omega$. В


этом случае в качестве $\wt H_n(\omega)$ можно принять


$\wt H_n(\omega)=H_n(\omega)P_n(\omega)$, откуда


$P_n(\omega)=H^{-1}_n(\omega)\wt H_n(\omega)$.





Ввиду наличия взаимнооднозначного соответствия между элементами


пространств $X_n$ и $Z_n$ вместо (9) можно рассматривать другое


приближенное уравнение


\begin{equation}


L_n(\omega)z_n(\omega)\equiv z_n(\omega)-\lambda B_n(\omega)


z_n(\omega)=\wt H_n(\omega)f(\omega)


\end{equation}


с непрерывным линейным оператором $B_n(\omega)$, действующим в


пространстве $Z_n$.


При этом $B_n(\omega)=H_n(\omega)F_n(\omega)H^{-1}_n(\omega)$,


$L_n(\omega)=H_n(\omega)K_n(\omega)H^{-1}_n(\omega)$.





Если найдено точное решение $\up{*}{z}_n(\omega)\in Z_n$


уравнения (10), то в


качестве приближенного решения уравнения (8) следует принять


$\ov x_n(\omega)=H^{-1}_n(\omega)\up{*}{z}_n(\omega)$,


где $\ov x_n(\omega)\in X_n$, $\omega\in\Omega$. В случае, когда


$Z_n=X_n$,


т.\,е. $H_n(\omega)$ рассматривается в качестве единичного оператора,


уравнения (9) и (10) а также их решения совпадают.





В области $\Omega\subset R^M$ зададим совокупность точек (узлов)


$\{\omega_{(i)}\}$,


где $\omega_{(i)}=(\omega_{(i)1},\dotsc,\omega_{(i)M})$,


$i=1,2,\dotsc,N$,


$N=N(n)$, $N\in\{1,2,\dotsc\}$. Поставим


в соответствие набору $\{\omega_{(i)}\}$ семейство


$\{\up{*}{z}_n(\omega_{(i)})\}$ частных


решений уравнения (10) так, что


$$


L_n(\omega_{(i)})\up{*}{z}_n(\omega_{(i)})\equiv


\up{*}{z}_n(\omega_{(i)})-\lambda B_n(\omega_{(i)})


\up{*}{z}_n(\omega_{(i)})=


\wt H_n(\omega_{(i)})f(\omega_{(i)}).


$$


Данные решения (которые по аналогии с [1], [2], [3] будем называть


опорными) строятся заранее с использованием высокоточных численных


методов и могут храниться в памяти ЭВМ.





Требуется на базе данного семейства построить приближенное


параметризованное решение $\wh z_n(\omega)\in Z_n$ уравнения (10),


которое


можно было бы использовать для нахождения приближенного решения


$\wh x_n(\omega)=H^{-1}_n(\omega)\wh z_n(\omega)$


исходного операторного


уравнения (8). При


этом должно выполняться неравенство


$$


\|\up{*}{x}(\omega)-\wh x_n(\omega)\|=\|\up{*}{x}(\omega)-


H^{-1}_n(\omega)\wh z_n(\omega)\|\le\delta,\quad


\wh x_n\in X,\ \; \omega\in\Omega,


$$


где $\up{*}{x}(\omega)$ --- точное параметризованное решение уравнения


(8), $\delta$


--- положительная постоянная.





Очевидно, количество и порядок расположения указанных выше


узлов $\{\omega_{(i)}\}$ зависят от типа используемых нормированных


пространств,


выбора области $\Omega\subset R^M$, свойств точного и приближенного


уравнений, а также требуемой точности построения приближенного


параметризованного решения уравнения (8).





По аналогии с (3), (4) и (5) пространства $X$, $Z_n$ и операторы


$F(\omega)$,


$B_n(\omega)$ для всех $\omega\in\Omega$ свяжем следующими условиями:


\begin{itemize}


\item[1)] для любого $x_n(\omega)\in X_n$


\begin{equation}


\|B_n(\omega)H_n(\omega)x_n(\omega)-\wt H_n(\omega)


F(\omega)x_n(\omega)\|\le\wt\eta_n\|x_n(\omega)\|,


\end{equation}


\item[2)] для любого $x(\omega)\in X$ найдется элемент


$x_n(\omega)\in X$


такой,


что


\begin{equation}


\|F(\omega)x(\omega)-x_n(\omega)\|\le\varepsilon_n\|x(\omega)\|,


\end{equation}


\item[3)] при некотором $f_n(\omega)\in X_n$


\begin{equation}


\|f(\omega)-f_n(\omega)\|\le\gamma_n\|f(\omega)\|,


\end{equation}


\end{itemize}


где константа $\gamma_n$, в отличие от констант $\wt\eta_n$ и


$\varepsilon_n$, в общем случае


зависит от $f(\omega)$.





Применительно же к задаче (8), (9), когда в качестве приближенного


решения уравнения (8) рассматривается точное решение уравнения (9),


вместо условия (11) используется новое условие: для любого


$x_n(\omega)\in X_n$


имеем


$\|P_n(\omega)F(\omega)x_n(\omega)-F_n(\omega)x_n(\omega)\|\le


\eta_n\|x_n(\omega)\|$,


где $\eta_n=\wt\eta_n\|H^{-1}_n(\omega)\|$. Условия (12), (13) не


затрагивают


операторов $F_n(\omega)$ и $P_n(\omega)$,


поэтому применительно к уравнению (9) не изменяются.


\medskip





{\bf 3. Построение параметризованных решений.} Допустим, что


для семейства опорных решений $\{\up{*}{z}_n(\omega_{(i)})\}$ (где


$\up{*}{z}_n(\omega_{(i)})\in Z_n$)


справедливы представления


\begin{gather}


\up{*}{z}_n(\omega_{(i)})=\varphi_n(\omega_{(i)},u_{(i)}),\quad


\omega_{(i)}\in\Omega\subset R^M, \notag \\


u_{(i)}\in U\subset R^K, \ \;i=1,2,\dotsc,N,\ \;


K=K(n),\ \;K\in\{1,2,\dotsc\},


\end{gather}


где $\varphi_n$ --- отображение, непрерывное в каждой из точек


$(\omega,u)\in\Omega\times U$,


$U$ --- ограниченное открытое связное множество;


$u_{(i)}=(u_{(i)1},\dotsc,u_{(i)K})$ ---


вектор с вещественными коэффициентами.





Для фиксированного $k\in\{1,2,\dotsc,K\}$ поставим в соответствие


семейству


узлов $\omega_{(1)},\dotsc,\omega_{(N)}$ набор коэффициентов


$u_{(1)k},\dotsc,u_{(N)k}$.


Проведем интерполяцию данного набора, сопоставив ему


скалярную вещественную функцию $\wh u_k(\omega)$ известного класса со


значениями


в $R^1$, непрерывную по $\omega\in\Omega$ [8], [9], [10]:


\begin{equation}


\wh u_k(\omega)=\mu(\omega,v_k),\quad


v_k\in V_k\subset R^N,\ \;\omega\in\Omega\subset R^M,


\end{equation}


где компоненты вектора $v_k=(v_{(1)k},\dotsc,v_{(N)k})$


находятся путем решения


следующей системы (как правило, линейной) алгебраических уравнений:


\begin{equation}


\wh u_k(\omega_{(i)})=\mu(\omega_{(i)},v_k)=u_{(i)k},\quad


i=1,2,\dotsc,N,\ \; k\in\{1,2,\dotsc,K\}.


\end{equation}


Здесь и далее полагаем, что соответствующие интерполяционные


задачи разрешимы и имеют единственное решение.


Теперь по аналогии с (15), (16) проводится интерполяция для


всех $k=1,2,\dotsc,K$, т.\,е. находится совокупность параметризованных


коэффициентов $\{\wh u_k(\omega)\}$, непрерывных по параметру


$\omega\in\Omega$. С


учетом этого приближенное параметризованное решение уравнения


(10), справедливое для всех $\omega\in\Omega$, может быть представлено


в виде


\begin{equation}


\wh z_n(\omega)=\varphi_n(\omega,\wh u(\omega)),\quad


\wh u(\omega)\in U,\ \;\wh z_n(\omega)\in Z_n,


\end{equation}


где $\wh u(\omega)=(\wh u_1(\omega),\dotsc,\wh u_K(\omega))$


--- векторная функция со значениями в $R^K$


непрерывная в $\Omega$.





Очевидно выполнение следующих характеристических свойств:


$$


\wh u(\omega_{(i)})=u_{(i)},\quad


\wh z_n(\omega_{(i)})=\varphi_n(\omega_{(i)},\wh u_{(i)})=


\varphi_n(\omega_{(i)},u_{(i)})=\up{*}{z}_n(\omega_{(i)}),\quad


i=1,2,\dotsc,N.


$$








Данные свойства показывают, что построенное параметризованное


решение (17) совпадает с опорными решениями семейства


$\{\up{*}{z}_n(\omega_{(i)})\}$


при фиксированных узлах интерполяции


$\{\omega_{(i)}\}$, где $\omega_{(i)}\in\Omega$, $i=1,2,\dotsc,N$.





По аналогии с (17) для точного решения уравнения (10) можно


воспользоваться представлением


\begin{equation}


\up{*}{z}_n(\omega)=\varphi_n(\omega,u(\omega)),\quad


u(\omega)\in U,\ \;\up{*}{z}_n(\omega)\in Z_n,


\end{equation}


где $u(\omega)=(u_1(\omega),\dotsc,u_K(\omega))$, при этом


$u(\omega_{(i)})=u_{(i)}$, $i=1,2,\dotsc,N$.





Если построено приближенное решение $\wh z_n(\omega)$


уравнения (10), то


в качестве приближенного решения $\wh x_n(\omega)$


уравнения (8) следует


принять $\wh x_n(\omega)=H^{-1}_n(\omega)\wh z_n(\omega)$, где


$\wh x_n(\omega)\in X_n$. Оценка


близости


приближенного решения $\wh x_n(\omega)$ уравнения (8) к точному решению


$\up{*}{x}(\omega)$


устанавливается в следующей теореме.





\begin{thm} Если выполнены условия $(11)$, $(12)$ и $(13)$, существует


непрерывный оператор $L^{-1}_n(\omega)$ и уравнение (8)


имеет непрерывное


в области $\Omega$ точное решение $\up{*}{x}(\omega)$, то и приближенное


решение


$\wh x_n(\omega)$ также непрерывно в области $\Omega$, и для всех


$\omega\in\Omega$ справедлива оценка


\begin{equation}


\|\up{*}{x}(\omega)-\wh x_n(\omega)\|=


\|\up{*}{x}(\omega)-H^{-1}_n(\omega)\wh z_n(\omega)\|\le


\wt p_n\|\up{*}{x}(\omega)\|+\|H^{-1}_n(\omega)\|\,


\|\Delta\varphi_n(\omega)\|,


\end{equation}


где


\begin{equation}


\Delta\varphi_n(\omega)=\varphi_n(\omega,u(\omega))-


\varphi_n(\omega,\wh u(\omega)),\quad


\Delta\varphi_n(\omega)\in Z_n,


\end{equation}


\begin{multline}


\wt p_n=[1+|\lambda|\varepsilon_n+\gamma_n\|K(\omega)\|]|\lambda|


\wt\eta_n\|H^{-1}_n(\omega)L^{-1}_n(\omega)\|+ \\


%


+[|\lambda|\varepsilon_n+\gamma_n\|K(\omega)\|][1+


\|H^{-1}_n(\omega)L^{-1}_n(\omega)\wt H_n(\omega)K(\omega)\|].


\end{multline}


\end{thm}





\begin{pf} В силу условий теоремы согласно


[4] для фиксированного $\omega\in\Omega$ справедлива оценка


\begin{equation}


\|\up{*}{x}(\omega)-\ov x_n(\omega)\|=


\|\up{*}{x}(\omega)-H^{-1}_n(\omega)\up{*}{z}_n(\omega)\|\le


\wt p_n\|\up{*}{x}(\omega)\|,


\end{equation}


где константа $\wt p_n$ определяется в соответствии с (21), а функция


$\ov x_n(\omega)$ является непрерывной в области $\Omega$.


Далее, с учетом неравенства треугольника


$$


\|\up{*}{x}(\omega)-\wh x_n(\omega)\|\le


\|\up{*}{x}(\omega)-\ov x_n(\omega)\|+


\|\ov x_n(\omega)-\wh x_n(\omega)\|,


$$


откуда, принимая во внимание (17), (18) и (22), получаем


\begin{multline}


\|\up{*}{x}(\omega)-\wh x_n(\omega)\|=


\|\up{*}{x}(\omega)-H^{-1}_n(\omega)\wh z_n(\omega)\|\le \\


%


\le \wt p_n\|\up{*}{x}(\omega)\|+


\|H^{-1}_n(\omega)\varphi_n(\omega,u(\omega))-


H^{-1}_n(\omega)\varphi_n(\omega,\wh u(\omega))\|.


\end{multline}





Так как линейный непрерывный оператор $H_n(\omega)$, действующий из


банахова пространства $X_n$ взаимно однозначно на банахово пространство


$Z_n$,


имеет непрерывный (а следовательно, ограниченный)


обратный оператор $H^{-1}_n(\omega)$, то с учетом (20) можно


воспользоваться


очевидным неравенством


\begin{multline*}


\|H^{-1}_n(\omega)\varphi_n(\omega,u(\omega))-


H^{-1}_n(\omega)\varphi_n(\omega,\wh u(\omega))\|\le \\


%


\le \|H^{-1}_n(\omega)\|\,\|\varphi_n(\omega,u(\omega))-


\varphi_n(\omega,\wh u(\omega))\|=


\|H^{-1}_n(\omega)\|\,\|\Delta\varphi_n(\omega)\|,\quad


\omega\in\Omega.


\end{multline*}


Отсюда с учетом (23) приходим к оценке (19).


\end{pf}





На основании данной теоремы легко получить оценку приближенного и


точного решений, не использующую данных, относящихся к


точному решению $\up{*}{x}(\omega)\in X$


$$


\|\up{*}{x}(\omega)-\wh x_n(\omega)\|\le


\frac{\wt p_n}{1-\wt p_n}\|\wh x_n(\omega)\|+


\|H^{-1}_n(\omega)\|\,\|\Delta\varphi_n(\omega)\|,\quad


\wt p<1,\ \;\omega\in\Omega.


$$


При заданном $N$ и оптимальном расположении узлов интерполяции


$\{\omega_{(i)}\}$


в области $\Omega$ (по аналогии с [1], [2], [3]) величина


$\|\up{*}{x}(\omega)-\wh x_n(\omega)\|$ достигает наименьшего значения


$$


\sup_\omega\inf_{\{\omega_{(i)}\}}


\|\up{*}{x}(\omega)-\wh x_n(\omega)\|\le


\sup_\omega\inf_{\{\omega_{(i)}\}}


[\wt p_n\|\up{*}{x}(\omega)\|+\|H^{-1}_n(\omega)\|\,


\|\Delta\varphi_n(\omega)\|]\le \delta.


$$





Условия сходимости последовательности


$\{\wh x_n(\omega)\}$ к точному решению $\up{*}{x}(\omega)$


уравнения (8) устанавливает





\begin{thm} Если оператор $K(\omega)$ имеет непрерывный обратный, и


при каждом $n=1,2,\dotsc$ выполнены условия $(11)$, $(12)$ и $(13)$,


причем


$$


\lim\wt\eta_n\|H_n(\omega)\|=\lim\varepsilon_n


\|H^{-1}_n(\omega)\wt H_n(\omega)\|=


\lim\gamma_n\|H^{-1}_n(\omega)\wt H_n(\omega)\|=0


$$


при $n\to\infty$ и


$\lim\|H^{-1}_n(\omega)\|\,\|\Delta\varphi_n(\omega)\|=0$ при


$n,N,K\to\infty$,


то при достаточно больших $n$ для всех $\omega\in\Omega$ имеет место


сходимость


последовательности решений


$\{\wh x_n(\omega)\}$ к точному решению $\up{*}{x}(\omega)$


и справедлива оценка


\begin{multline*}


\|\up{*}{x}(\omega)-\wh x_n(\omega)\|\le c_{n0}\wt\eta_n


\|H^{-1}_n(\omega)\|+c_{n1}\varepsilon_n


\|H^{-1}_n(\omega)\wt H_n(\omega)\|+ \\


%


+c_{n2}\gamma_n\|H^{-1}_n(\omega)\wt H_n(\omega)\|+


\|H^{-1}_n(\omega)\|\,\|\Delta\varphi_n(\omega)\|.


\end{multline*}


\end{thm}





Доказательство теоремы 2 непосредственно следует из неравенства


треугольника и соответствующих оценок, приведенных в


[4]. Соотношение


$\lim\|H^{-1}_n(\omega)\|\,\|\Delta\varphi_n(\omega)\|=0$ при


$n,N,K\to\infty$ указывает на то, что методы


интерполяции, используемые при построении


параметризованного решения $\wh x_n(\omega)$, являются сходящимися.





\begin{note} При построении параметризованных решений операторных


уравнений описанным выше способом вместо рассмотренных выше


процедур интерполяции могут использоваться известные процедуры


аппроксимации (приближения). Очевидно, что в этом случае отмеченные


ранее характеристические свойства могут не выполняться, при


этом в качестве $\|\Delta\varphi_n(\omega)\|$ рассматриваются нормы


соответствующих


остаточных членов аппроксимации.


\end{note}





\begin{note} Если рассмотренный выше подход к построению


параметризованных решений применяется непосредственно к уравнению


(9), то необходимо поставить в соответствие набору узлов


$\{\omega_{(i)}\}$


семейство $\{\up{*}{x}_n(\omega_{(i)})\}$


опорных решений этого уравнения. По аналогии


с (14)--(18) точное и приближенное решения уравнения (9) можно


представить в виде


$\up{*}{x}_n(\omega)=\xi_n(\omega,\psi(\omega))$,


$\wh x_n(\omega)=\xi_n(\omega,\wh\psi(\omega))$,


где


$\wh\psi(\omega)=(\wh\psi_1(\omega),\dotsc,\wh\psi_K(\omega))$,


$\wh\psi_k(\omega)=\theta(\omega,w_k)$, $w_k\in{}W_k\subset R^N$,


$k\in\{1,2,\dotsc,K\}$.


Здесь компоненты вектора $w_k=(w_{(1)k},\dotsc,w_{(N)k})$ находятся из


системы уравнений


$\wh\psi_k(\omega_{(i)})=\theta(\omega_{(i)},w_k)=\psi_{(i)k}$,


$i=1,2,\dotsc,N$, $k\in\{1,2,\dotsc,K\}$.


Для оценки близости


$\up{*}{x}(\omega)$ и $\wh x_n(\omega)$ можно воспользоваться оценкой


$\|\up{*}{x}(\omega)-\wh x_n(\omega)\|\le p_n


\|\up{*}{x}(\omega)\|+\|\Delta\xi_n(\omega)\|$,


где


$\Delta\xi_n(\omega)=\xi_n(\omega,\psi(\omega))-


\xi_n(\omega,\wh\psi(\omega))$,


$p_n=2|\lambda|\eta_n\|K^{-1}_n(\omega)\|+


[|\lambda|\varepsilon_n+\gamma_n\|K(\omega)\|]


[1+\|K^{-1}_n(\omega)P_n(\omega)K(\omega)\|]$.


\end{note}





{\bf 4. Параметризованные решения метода Галеркина.} Остановимся


подробнее на важном для практики случае, когда $Z_n=X_n$, а в


качестве $F_n(\omega)$ рассматривается оператор $P_n(\omega)F(\omega)$,


что позволяет


перейти от (9) к новому приближенному уравнению


\begin{equation}


K_n(\omega)x_n(\omega)\equiv x_n(\omega)-


\lambda P_n(\omega)F(\omega)x_n(\omega)=


P_n(\omega)f(\omega).


\end{equation}


Очевидно, что в этом случае можно положить $\wt\eta_n=0$, и,


следовательно,


формулировки приведенных выше теорем упрощаются.





Пусть $X_n$ --- конечномерное (размерности $n$) пространство $X$.


Каждый элемент $x_n(\omega)\in X_n$ с учетом второго замечания по


аналогии


с (18) единственным образом представим в форме (при этом $K=n$)


\begin{equation}


x_n(\omega)=\sum^n_{k=1}\psi_k(\omega)\xi_k,\quad \xi_k\in X_n,


\end{equation}


где $\psi(\omega)=(\psi_1(\omega),\dotsc,\psi_n(\omega))$ --- вектор


искомых параметризованных коэффициентов,


а совокупность $\{\xi_k\}$ образует базис в $X_n$.





Рассмотрим полную в $X_n$ систему $\{Q_j(\omega)\}$ линейных


параметризованных


функционалов $Q_j(\omega)$ таких, что из равенств


$Q_j(\omega)x_n(\omega)=0$, $j=1,2,\dotsc,n$,


следует $x_n(\omega)=0$. В этом случае вместо (24) можно


ограничиться рассмотрением системы равенств


$$


Q_j(\omega)P_n(\omega)K(\omega)x_n(\omega)=


Q_j(\omega)P_n(\omega)f(\omega),\quad j=1,2,\dotsc,n.


$$


Разыскивая решение уравнения (24) в форме (25), приходим к


параметризованной системе линейных алгебраических уравнений метода


Галеркина а абстрактной форме (по аналогии с [4], [5])


\begin{equation}


\sum^n_{k=1}\psi_k(\omega)Q_j(\omega)\xi_k-\lambda


\sum^n_{k=1}\psi_k(\omega)Q_j(\omega)P_n(\omega)


F(\omega)\xi_k=Q_j(\omega)P_n(\omega)f(\omega),\quad


j=1,2,\dotsc,n.


\end{equation}





Если система функционалов $\{Q_j(\omega)\}$ биортогональна базису


$\{\xi_k\}$, то с учетом (26) получаем


$$


\psi_j(\omega)-\lambda\sum^n_{k=1}\psi_k(\omega)


Q_j(\omega)P_n(\omega)F(\omega)\xi_k=


Q_j(\omega)P_n(\omega)f(\omega),\quad j=1,2,\dotsc,n.


$$


В частности, если $X$ --- гильбертово пространство, а $P_n(\omega)$ ---


оператор ортогонального проектирования, то, предполагая систему


$\{\xi_k\}$ ортонормальной, имеем


\begin{equation}


\psi_j(\omega)-\lambda\sum^n_{k=1}\psi_k(\omega)


\la F(\omega)\xi_k,\xi_j\ra=


\la f(\omega),\xi_j\ra,\quad j=1,2,\dotsc,n,


\end{equation}


где $\la\cdot,\cdot\ra$ --- символ скалярного произведения.





Решая систему уравнений (27) относительно $\{\psi_k(\omega)\}$ для


совокупности


узлов $\{\omega_{(i)}\}$ (где $i=1,2,\dotsc,N$),


найдем семейство опорных


решений $\{\up{*}{x}_n(\omega_{(i)})\}$


уравнения (24), которое с учетом второго замечания


по аналогии с (14) может быть представлено в виде


$\{\up{*}{x}_n(\omega_{(i)})=\xi\psi^T_{(i)}\}$,


где $\xi=(\xi_1,\dotsc,\xi_n)$ --- вектор базисных функций,


$\psi_{(i)}=(\psi_{(i)1},\dotsc,\psi_{(i)n})$ --- вектор


искомых коэффициентов, которые находятся из (27) для узла


$\omega_{(i)}\in\Omega$,


при этом $\psi_{(i)k}=\psi_k(\omega_{(i)})$, $k=1,2,\dotsc,n$; $T$ ---


символ транспонирования.





По аналогии с (15), (16) с учетом второго замечания можно


воспользоваться следующей интерполяцией:


\begin{equation}


\wh\psi_k(\omega)=\sum^N_{j=1}w_{(j)k}\theta_{(j)}(\omega),


\quad k=1,2,\dotsc,n,


\end{equation}


где совокупность коэффициентов $w_{(1)k},\dotsc,w_{(N)k}$ для


фиксированного $k$


находится как решение системы уравнений


\begin{equation}


\wh\psi_k(\omega_{(i)})=\sum^N_{j=1}


w_{(i)k}\theta_{(j)}(\omega_{(i)})=


\psi_{(i)k},\quad i=1,2,\dotsc,N.


\end{equation}





Приближенное параметризованное решение уравнения (24) может


быть представлено в виде


$\wh x_n(\omega)=\xi\wh\psi\,{}^T(\omega)$, где


$\wh\psi(\omega)=(\wh\psi_1(\omega),\dotsc,\wh\psi_n(\omega))$.


При этом


$\wh x_n(\omega_{(i)})=\xi\wh\psi\,{}^T(\omega_{(i)})=


\xi\psi^T_{(i)}=\up{*}{x}_n(\omega_{(i)})$, $i=1,2,\dotsc,N$.


Если в качестве точного решения уравнения (24) принять


$\up{*}{x}_n(\omega)=\xi\psi^T(\omega)$,


то с учетом второго замечания


$\|\up{*}{x}(\omega)-\wh x_n(\omega)\|\le


p_n\|\up{*}{x}(\omega)\|+\|\xi\Delta\psi^T(\omega)\|$,


где


$\Delta\psi(\omega)=\psi(\omega)-\wh\psi(\omega)$,


$p_n=[|\lambda|\varepsilon_n+\gamma_n\|K(\omega)\|]


[1+\|K^{-1}_n(\omega)P_n(\omega)K(\omega)\|]$.





Можно воспользоваться несколько иным подходом к построению


параметризованного решения $\wh x_n(\omega)$


по сравнению с рассмотренным


выше. Для этого представим $\wh x_n(\omega)$ в следующем виде


\begin{equation}


\wh x_n(\omega)=\sum^N_{j=1}\sum^n_{k=1}


w_{(j)k}\theta_{(j)}(\omega)\xi_k.


\end{equation}


Задавшись совокупностью узлов $\{\omega_{(i)}\}$,


составим с учетом (27)


и (30) систему линейных алгебраических уравнений относительно


массива $\{w_{(j)k}\}$ искомых коэффициентов $w_{(j)k}$


\begin{multline}


\sum^N_{p=1}w_{(p)j}\theta_{(p)}(\omega_{(i)})-


\lambda\sum^n_{k=1}\sum^N_{d=1}w_{(d)k}\theta_{(d)}


(\omega_{(i)})\la F(\omega_{(i)})\xi_k,\xi_j\ra = \\


%


=\la f(\omega_{(i)}),\xi_j\ra,\quad j=1,2,\dotsc,n,\ \;


i=1,2,\dotsc,N.


\end{multline}


Предполагается, что при выбранной совокупности узлов


$\{\omega_{(i)}\}$ система (31) не вырождена.


\medskip





{\bf 5. Пример.} По аналогии с [4] рассмотрим краевую задачу для


обыкновенного дифференциального уравнения


\begin{equation}


x^{(2P)}(\omega,t)-\lambda\sum^{2P}_{r=1}p_r(\omega,t)


x^{(2P-r)}(\omega,t)=y(\omega,t),\quad


t\in[a,b]\subset R^1,\ \;\omega\in[c,d]\subset R^1


\end{equation}


при однородных краевых условиях


\begin{equation}


x^{(r)}(\omega,a)=x^{(r)}(\omega,b)=0,\quad


r=0,1,\dotsc,P-1,


\end{equation}


где для фиксированного $\omega\in\Omega$ выполняется следующее:


$x(\omega,t)\in X=C^{(2P)}[a,b]$,


$p_r(\omega,t)\in C^{(1)}[a,b]$, $y(\omega,t)\in Y=C^{(1)}[a,b]$.


Будем рассматривать уравнение


(32) как функциональное в пространстве $2P$-непрерывно дифференцируемых


по $t$ функций $x(\omega,t)$, удовлетворяющих краевым условиям


(33). В этом случае с учетом (1) имеем $x(\omega)=x(\omega,t)$,


$G(\omega)=G$, $Gx(\omega,t)=x^{(2P)}(\omega,t)$,


$$


T(\omega)x(\omega,t)=\sum^{2P}_{r=1}p_r(\omega,t)


x^{(2P-r)}(\omega,t).


$$


Полагаем также, что выполняются условия теоремы о неявной


функции, обеспечивающие требуемую гладкость решения


$\up{*}{x}(\omega,t)$ задачи


(32), (33) по параметру $\omega$ в замкнутом шаре $\Omega=[c,d]$.





За $X_n$ примем (с учетом (25)) множество функций вида


\begin{equation}


x_n(\omega,t)=(t-a)^P(b-t)^P\sum^n_{k=1}


\psi_k(\omega)t^{k-1}=\sum^n_{k=1}\psi_k(\omega)


\xi_k(t),


\end{equation}


так что краевые условия (33) удовлетворяются автоматически.





Если воспользоваться интерполяционным методом (иначе методом


совпадения (коллокации) [4]), то коэффициенты


$\psi_{(i)k}=\psi_k(\omega_{(i)})$, $k=1,2,\dotsc,n$,


соответствующие $i$-му опорному решению


$\up{*}{x}_n(\omega_{(i)},t)$,


могут быть найдены из условия удовлетворения уравнения (32) в


некоторой заданной системе точек (узлов) коллокации


$t_1,t_2,\dotsc,t_n$


\begin{equation}


x^{(2P)}_n(\omega_{(i)},t_j)-\lambda\sum^{2P}_{r=1}


p_r(\omega_{(i)},t_j)x^{(2P-r)}(\omega_{(i)},t_j)=


y(\omega_{(i)}t_j),\quad j=1,2,\dotsc,n.


\end{equation}





Решая систему (35) для каждого $i\in\{1,2,\dotsc,N\}$, поставим в


соответствие набору $\{\omega_{(i)}\}$ семейство


$\{\up{*}{x}_n(\omega_{(i)},t)\}$ частных решений


уравнения (2), которое в нашем случае принимает следующий вид:


$Gx_n(\omega)-\lambda\Phi_n Tx_n(\omega)=\Phi_n y(\omega)$,


где $x_n(\omega)=x_n(\omega,t)$, $x_n(\omega)=\wt x(\omega)$,


$\Phi_n=\wt\Phi(\omega)$.


Отображение $\Phi_n$ пространства $Y$ на $\wt Y=Y_n$ определим,


принимая в качестве $y_n(\omega,t)=\Phi_n y(\omega,t)$


для фиксированного


$\omega$ интерполяционный


полином $(n-1)$-й степени, имеющий в узлах $t_1,t_2,\dotsc,t_n$


те же значения, что и $y(\omega,t)$.





По аналогии с [1]--[3], используя интерполяцию на основе многочлена


Лагранжа с учетом (28), (29) имеем


\begin{equation}


\wh\psi_k(\omega)=\sum^N_{j=1}w_{(j)k}\theta_{(j)}(\omega)=


\sum^N_{j=1}\psi_{(j)k}L_{(j)}(\omega),


\end{equation}


где


\begin{gather*}


w_{(j)k}=\psi_{(j)k},\quad \theta_{(j)}(\omega)=


L_{(j)}(\omega), \\


L_{(j)}(\omega)=\prod^N\begin{Sb} s=1 \\ s\ne j\end{Sb}


\frac{\omega-\omega_{(s)}}{\omega_{(j)}-\omega_{(s)}},\quad


\omega_{(j)},\omega_{(s)}\in\Omega=[c,d], \\


\wh\psi_k(\omega_{(i)})=\sum^N_{j=1}\psi_{(j)k}L_{(j)}


(\omega_{(i)})=\psi_{(j)k},\quad i=1,2,\dotsc,N,\ \;


k=1,2,\dotsc,n.


\end{gather*}





Таким образом, в силу (34), (36) приближенное параметризованное


решение краевой задачи (32), (33) имеет вид


\begin{equation}


\wh x_n(\omega,t)=(t-a)^P(b-t)^P


\sum^n_{k=1}t^{k-1}\sum^N_{j=1}\psi_{(j)k}\prod^N_{s=1}


\frac{\omega-\omega_{(s)}}{\omega_{(j)}-\omega_{(s)}},


\end{equation}


при этом


$\wh x_n(\omega_{(i)},t_r)=\up{*}{x}_n(\omega_{(i)},t_r)$,


$i=1,2,\dotsc,N$, $r=1,2,\dotsc,n$.





Применительно к данному примеру условия (3), (4) и (5) выполнены


соответственно с константами $\mu_0=0$, $\mu_1=O(1/n)$ и


$\mu_2=O(1/n)$, если


вспомнить, что по условию $y(\omega,t)$ --- дважды непрерывно


дифференцируемая


(по $t$) функция. В случае, когда узлы $\{t_r\}$ чебышевские,


справедливо


$\|\up{*}{x}(\omega,t)-\up{*}{x}_n(\omega,t)\|=O(\ln n/n)$.





По аналогии с [2], [3] для оценки погрешности интерполяции можно


воспользоваться неравенством


\begin{equation}


|\psi_k(\omega)-\wh\psi_k(\omega)|\le\frac{Q_N}{N\,!}


\prod^N_{i=1}|\omega-\omega_{(i)}|,\quad \omega\in[c,d],


\end{equation}


при этом узлы $\{\omega_{(i)}\}$ являются чебышевскими, т.\,е.


$$


\omega_{(i)}=\frac{d+c}{2}-\frac{d-c}{2}\cos


\biggl(\frac{2i-1}{2N}\pi\biggr),


$$


то


$$


\max_{\omega\in[c,d]}\prod^N_{i=1}|\omega-\omega_{(i)}|=


\frac{(d-c)^N}{2^{2N-1}}.


$$


Таким образом, с учетом (34), (37) и (38) имеем оценку


$$


|\up{*}{x}_n(\omega,t)-\wh x_n(\omega,t)|\le


|(t-a)^P(b-t)^P|\frac{Q_N}{N\,!}\frac{(d-c)^N}{2^{2N-1}}


\sum^n_{k=1}|t|^{k-1}


$$


при этом


$|\up{*}{x}_n(\omega,t)-\wh x_n(\omega,t)|\to 0$ при $n,N\to\infty$ для


всех $\omega\in[c,d]$


и $t\in[a,b]$. Таким образом, равномерная сходимость


$$


\max_{t,\omega}|\up{*}{x}(\omega,t)-\wh x_n(\omega,t)|\to 0,\quad


t\in[a,b],\ \; \omega\in[c,d],


$$


при $n,N\to\infty$, с учетом свойств гладкости приближенного и точного


решений задачи (32), (33), следует из сходимости метода коллокации


и сходимости метода интерполяции многочлена Лангража.





\begin{thebibliography}{99}





\bibitem{1}


Булычев Ю.Г. {\em Метод опорных интегральных кривых решения задачи


Коши для обыкновенных дифференциальных уравнений} // Журн. вычисл.


матем. и матем. физ. -- 1988. -- Т.\,28. -- \No\,10.


-- С.\,1482--1490.


\bibitem{2}


Булычев Ю.Г. {\em Методы численно-аналитического интегрирования


дифференциальных уравнений} // Журн. вычисл. матем. и матем. физ. --


1991. -- Т. 31. -- \No\,9. -- С.\,1305--1319.


\bibitem{3}


Булычев Ю.Г. {\em Численно-аналитическое интегрирование


дифференциальных уравнений с использованием обобщенной интерполяции}


// Журн. вычисл. матем. и матем. физ.


-- 1994. -- Т.\,34. -- \No\,4. -- С.\,520--532.


\bibitem{4}


Канторович Л.В., Акилов Г.П. {\em Функциональный анализ}. -- 3-е изд. --


М.: Наука, 1984. -- 752 с.


\bibitem{5}


Красносельский М.А., Вайникко Г.М., Забрейко П.П.


{\em Приближенное решение операторных


уравнений}. -- М.: Наука, 1969. -- 455 с.


\bibitem{6}


Вайникко Г.М. {\em Возмущенный метод Галеркина и общая теория


приближенных методов для нелинейных уравнений} // Журн. вычисл. матем.


и матем. физ. -- 1967. -- Т.\,7. -- \No\,4. -- С.\,723--751.


\bibitem{7}


Михлин С.Г. {\em Вариационные методы в математической физике}. --


2-е изд., перераб. и доп. -- М.: Наука, 1970. -- 512 с.


\bibitem{8}


Березин И.С., Жидков Н.П. {\em Методы вычислений}. Т.\,1. --


3-е изд., перераб. и доп. -- М.: Наука, 1966. -- 632 с.


\bibitem{9}


Бабенко К.И. {\em Основы численного анализа}. -- М.: Наука, 1986.


-- 744 с.


\bibitem{10}


Иванов В.В. {\em Методы вычислений на ЭВМ}: Справ. пособие. --


Киев: Наук. думка, 1986. -- 582~с.





\end{thebibliography}





\vskip 0.5cm





\postup{\it Ростовское высшее военное}{$\qquad$\it Поступили}


\postup{\it командно-инженерное}{{\it первый вариант} 12.02.1996}


\postup{\it училище ракетных войск}{{\it окончательный


вариант} 02.12.1996}


\label{end}


\end{document}


